
\section{Redirect and Callback}

\textit{Normative language per \href{https://tools.ietf.org/html/rfc2119}{RFC~2119}.}

\subsection{Callback URI}

The callback URI is tenant-specific. Each tenant deploys the Canvasflow SPA on 
their own domain, and the callback URI is formed from that domain with a 
standardized path:

```
https://www.<tenant-domain>/callback
```

For example, a tenant whose domain is `example.com` uses:

```
https://www.example.com/callback
```

This URI MUST be registered in the tenant's IdP exactly as shown --- including the HTTPS
scheme, the exact hostname, and the \texttt{/callback} path. No trailing slash. No
wildcards. No query parameters.



The tenant's IdP MUST validate the \texttt{redirect\_uri} parameter in every
authorization request against this registered value using exact string matching
(\href{https://tools.ietf.org/html/rfc6749#section-3.1.2}{RFC~6749~\S3.1.2}). 
Partial matching, wildcard matching, or prefix matching MUST NOT be permitted. 
Exact-match validation prevents open-redirect attacks, where an attacker substitutes 
a malicious redirect URI to capture the authorization code.

The tenant MUST NOT register additional redirect URIs for the Canvasflow SPA client
registration beyond this exact value, unless required for a sandbox or staging
environment (in which case, a separate client registration SHOULD be used for each
environment).

\subsection{Authorization Code Flow with PKCE}\label{sec:authorization-with-pkce}

This integration uses the OAuth~2.0 Authorization Code flow 
(\href{https://tools.ietf.org/html/rfc6749#section-4.1}{RFC~6749~\S4.1}) combined
with Proof Key for Code Exchange 
(\href{https://tools.ietf.org/html/rfc7636}{RFC~7636}) with 
\texttt{response\_mode=fragment}.

\subsubsection{What Travels in the Fragment}

Only the short-lived, single-use \textbf{authorization code} travels in the URL fragment:


\begin{lstlisting}[basicstyle=\footnotesize\ttfamily]
https://www.<tenant-domain>/callback#code=<authorization_code>&state=<state>
\end{lstlisting}

Tokens --- access token, refresh token, and ID token --- are \textbf{never} placed in the
URL. They are obtained by the SPA in a direct HTTP POST to the tenant's token endpoint,
which returns them in the response body over a TLS-secured connection.

The fragment is cleared from the browser URL bar immediately after the SPA reads the
code, using \texttt{history.replaceState()}. The authorization code's exposure in the
fragment is bounded: it is single-use, it has a short server-side lifetime (typically
$\sim$60~seconds), and it is useless to an attacker without the \texttt{code\_verifier},
which never leaves the SPA.

\subsubsection{Retirement of Implicit-Style Token Fragment Delivery}

Prior versions of this integration specification described delivering tokens
(\texttt{access\_token}, \texttt{refresh\_token}, \texttt{id\_token}) directly in the
URL fragment --- the Implicit flow pattern (\href{https://tools.ietf.org/html/rfc6749#section-4.2}{RFC~6749~\S4.2}). This pattern is
\textbf{retired} in this specification and MUST NOT be used.

\textbf{Justification.} The Implicit flow is deprecated by 
\href{https://tools.ietf.org/html/rfc9700#section-2.1.2}{RFC~9700~\S2.1.2} for the
following reasons:

\begin{itemize}
  \item Tokens in URLs are recorded in browser history, where they persist until the user
    manually clears history.
  \item Browser extensions can read \texttt{window.location.hash} and exfiltrate tokens.
  \item If the callback page loads any third-party resource (analytics, fonts, CDN), the
    token may appear in the \texttt{Referer} header of outbound requests from the page.
  \item Access tokens have no single-use constraint --- an intercepted token is valid until
    \texttt{exp}, which may be up to 3600~seconds.
\end{itemize}

The Authorization Code + PKCE pattern (\hyperref[sec:authorization-with-pkce]{\S7.2}) 
eliminates these risks. The authorization code that travels in the fragment is 
single-use, short-lived, and provides no value to an attacker without the 
\texttt{code\_verifier}. Tokens are obtained via a direct,
TLS-protected HTTP POST that does not appear in any URL or HTTP log.

\subsection{PKCE Requirements}

PKCE MUST be implemented using the \texttt{S256} challenge method. The \texttt{plain}
method MUST NOT be used (\href{https://tools.ietf.org/html/rfc7636#section-4.2}{RFC~7636~\S4.2}).

\subsubsection{\texttt{code\_verifier} Generation}

The \texttt{code\_verifier} MUST be generated using a cryptographically secure random
source. In the browser, this MUST be the Web Crypto API (\texttt{crypto.getRandomValues}).
The \texttt{code\_verifier} MUST be between 43 and 128 characters in length, using the
character set \texttt{[A-Za-z0-9\texttt{-}.\texttt{\_}\textasciitilde]}
(unreserved characters per \href{https://tools.ietf.org/html/rfc7636#section-4.1}{RFC~7636~\S4.1}).

\begin{lstlisting}[language=js]
// Correct -- cryptographically secure
const array = new Uint8Array(32);
crypto.getRandomValues(array);
const codeVerifier = base64urlEncode(array); // 43 chars minimum from 32 bytes
\end{lstlisting}

The \texttt{code\_verifier} MUST be stored in memory only during the authorization flow
and discarded after the token exchange is complete. It MUST NOT be persisted to
\texttt{localStorage}, \texttt{sessionStorage}, or any other durable storage.

\subsubsection{\texttt{code\_challenge} Derivation}\label{sec:code-challenge-derivation}

The \texttt{code\_challenge} is derived from the \texttt{code\_verifier} as:

\begin{lstlisting}
code_challenge = BASE64URL(SHA256(ASCII(code_verifier)))
\end{lstlisting}

This derivation MUST also use the Web Crypto API:

\begin{lstlisting}[language=js]
const encoder = new TextEncoder();
const data = encoder.encode(codeVerifier);
const digest = await crypto.subtle.digest('SHA-256', data);
const codeChallenge = base64urlEncode(new Uint8Array(digest));
\end{lstlisting}

\subsubsection{Why PKCE Is Required Even When the IdP Allows Skipping It}

PKCE is a defense-in-depth requirement --- it is mandatory for this integration regardless
of whether the tenant's IdP enforces it. PKCE protects against two distinct attack
classes:

\begin{enumerate}
  \item \textbf{Authorization-code interception (\href{https://tools.ietf.org/html/rfc7636\#section-1}{RFC~7636~\S1}).} An attacker who
    intercepts the authorization code from the URL fragment cannot exchange it for tokens
    without the \texttt{code\_verifier}, which is never transmitted over any channel
    until the token exchange.

  \item \textbf{Authorization-code injection (\href{https://tools.ietf.org/html/rfc9700\#section-4.4}{RFC~9700~\S4.4}).} PKCE prevents an
    attacker from injecting a stolen authorization code (obtained from a separate session
    or system) into the victim's flow. The \texttt{code\_verifier} ties the token
    exchange to the specific browser session that initiated the authorization request.
\end{enumerate}

Tenants whose IdP marks PKCE as optional MUST configure the Canvasflow SPA client
registration to enforce PKCE.

\subsection{\texttt{state} Parameter --- CSRF Protection}

The \texttt{state} parameter MUST be included in every authorization request
(\href{https://tools.ietf.org/html/rfc6749\#section-4.1.1}{RFC~6749~\S4.1.1}). 
The SPA MUST:

\begin{enumerate}
  \item Generate a cryptographically random \texttt{state} value before redirecting to
    the IdP.
  \item Store the \texttt{state} value in memory (not \texttt{localStorage}).
  \item Verify that the \texttt{state} value returned in the callback fragment exactly
    matches the stored value before accepting the authorization code.
  \item Reject and discard the authorization code if the \texttt{state} values do not
    match.
\end{enumerate}

A mismatched or absent \texttt{state} indicates a potential Cross-Site Request Forgery
(CSRF) attack targeting the callback endpoint.

\subsection{\texttt{nonce} Parameter --- ID Token Replay Protection}



When the \texttt{openid} scope is requested (which produces an ID token), the
\texttt{nonce} parameter MUST be included in the authorization request
(\href{https://openid.net/specs/openid-connect-core-1_0.html#AuthRequest}{OIDC Core \S3.1.2.1}).

The \texttt{nonce} MUST be a cryptographically random value generated fresh for each
authorization request. The SPA MUST verify that the \texttt{nonce} claim in the received
ID token matches the value generated for that request. A mismatched \texttt{nonce}
indicates a potential ID token replay attack.

The \texttt{nonce} applies to the ID token only. The access token does not carry a
\texttt{nonce} claim.

\subsection{Token Endpoint Authentication}\label{sec:token-endpoint-authentication}

The Canvasflow SPA is a \textbf{public OAuth~2.0 client} 
(\href{https://tools.ietf.org/html/rfc6749\#section-2.1}{RFC~6749~\S2.1}). It MUST NOT
use a client secret. Token endpoint authentication method MUST be set to \texttt{none}
in the tenant's client registration.

The absence of a client secret is correct and intentional: client secrets in
browser-based applications cannot be kept confidential and provide no security benefit
(\href{https://tools.ietf.org/html/rfc9700\#section-2.1}{RFC~9700~\S2.1}). 
PKCE provides equivalent CSRF protection for public clients.

\subsection{Authorization Request Summary}

The following parameters MUST be included in every authorization request sent to the
tenant's IdP:

\begin{longtable}{@{}llp{6cm}@{}}
\toprule
\textbf{Parameter} & \textbf{Required} & \textbf{Value} \\
\midrule
\endhead
\texttt{response\_type}       & MUST           & \texttt{code} \\[4pt]
\texttt{client\_id}           & MUST           & Tenant-assigned client ID for the Canvasflow SPA \\[4pt]
\texttt{redirect\_uri}        & MUST           & \texttt{https://app.canvasflow.io/callback} \\[4pt]
\texttt{scope}                & MUST           & At minimum \texttt{openid}; additional scopes at tenant's discretion \\[4pt]
\texttt{code\_challenge}      & MUST           & Derived from \texttt{code\_verifier} per \hyperref[sec:code-challenge-derivation]{\S7.3.2} \\[4pt]
\texttt{code\_challenge\_method} & MUST        & \texttt{S256} \\[4pt]
\texttt{state}                & MUST           & Cryptographically random CSRF token \\[4pt]
\texttt{response\_mode}       & SHOULD         & \texttt{fragment} \\[4pt]
\texttt{nonce}                & MUST (when \texttt{openid} in scope) & Cryptographically random per-request value \\[4pt]
\texttt{audience}             & MUST (Auth0 only) & \texttt{canvasflow-api} \\
\bottomrule
\end{longtable}
